% Part: first-order-logic
% Chapter: tableaux
% Section: provability-consistency

\documentclass[../../../include/open-logic-section]{subfiles}

\begin{document}

\iftag{FOL}
      {\olfileid{fol}{tab}{prv}}
      {\olfileid{pl}{tab}{prv}}

\olsection{\usetoken{S}{derivability} والاتساق}

سنثبت الآن عددًا من خصائص علاقة ال!!{derivability}. وهذه الخصائص مهمة في
ذاتها، وسيكون لكل منها دور في برهان مبرهنة الاكتمال.

\begin{prop}\ollabel{prop:provability-contr}
  إذا كان $\Gamma \Proves !A$ وكانت $\Gamma \cup \{!A\}$ غير متسقة،
  فإن $\Gamma$ غير متسقة.
\end{prop}

\begin{proof}
توجد مجموعتان منتهيتان $\Gamma_0 = \{!B_1, \dots, !B_n\}$ و$\Gamma_1
  =\{!C_1, \dots, !C_m\} \subseteq \Gamma$ بحيث
  \begin{align*}
    \{\sFmla{\False}{!A}, &
    \sFmla{\True}{!B_1}, \dots, \sFmla{\True}{!B_n}\} \\
    \{\sFmla{\True}{!A}, &
    \sFmla{\True}{!C_1}, \dots, \sFmla{\True}{!C_m}\}
  \end{align*}
  لهما !!{tableau}s مغلقة. وباستخدام قاعدة \Cut{} على $!A$ يمكننا
  جمعهما في !!{tableau} مغلق واحد يثبت أن $\Gamma_0
  \cup \Gamma_1$ غير متسقة. وبما أن $\Gamma_0
  \subseteq \Gamma$ و$\Gamma_1 \subseteq \Gamma$، فإن $\Gamma_0 \cup
  \Gamma_1 \subseteq \Gamma$، ومن ثم
  فــ$\Gamma$ غير متسقة.
\end{proof}

\begin{prop}
\ollabel{prop:prov-incons}
$\Gamma \Proves !A$ إذا وفقط إذا كانت $\Gamma \cup \{\lnot !A\}$ غير متسقة.
\end{prop}

\begin{proof}
افترض أولًا أن $\Gamma \Proves !A$، أي أن هناك !!{tableau} مغلقًا لــ
\[
\{\sFmla{\False}{!A},
\sFmla{\True}{!B_1}, \dots, \sFmla{\True}{!B_n}\}
\]
وباستخدام قاعدة $\TRule{\True}{\lnot}$ يمكن تحويله إلى !!{tableau}
مغلق لــ
\[
\{\sFmla{\True}{\lnot !A},
\sFmla{\True}{!B_1}, \dots, \sFmla{\True}{!B_n}\}.
\]

ومن جهة أخرى، إذا كان للأخير !!{tableau} مغلق، فيمكننا تحويله إلى
!!{tableau} مغلق للأول بحذف كل صيغة تنتج من تطبيق
\TRule{\True}{\lnot} على الافتراض الأول~$\sFmla{\True}{\lnot !A}$،
وحذف ذلك الافتراض أيضًا، ثم إضافة الافتراض $\sFmla{\False}{!A}$. ذلك
أن الفرع إن كان مغلقًا من قبل لاحتوائه على نتيجة تطبيق
\TRule{\True}{\lnot} على $\sFmla{\True}{\lnot !A}$، أي
$\sFmla{\False}{!A}$، فإن الفرع الموافق له في ال!!{tableau} الجديد مغلق
أيضًا. وإذا كان فرع في الجدول الدلالي القديم مغلقًا لاحتوائه على الافتراض
$\sFmla{\True}{\lnot !A}$ وعلى $\sFmla{\False}{\lnot
  !A}$ معًا، فيمكننا
تحويله إلى فرع مغلق بتطبيق $\TRule{\False}{\lnot}$ على
$\sFmla{\False}{\lnot !A}$ لنحصل على $\sFmla{\True}{!A}$. وهذا يغلق
الفرع لأننا أضفنا $\sFmla{\False}{!A}$ بوصفه افتراضًا.
\end{proof}

\begin{prob}
أثبت أن $\Gamma \Proves \lnot !A$ إذا وفقط إذا كانت $\Gamma \cup \{!A\}$ غير متسقة.
\end{prob}

\begin{prop}\ollabel{prop:explicit-inc}
  إذا كان $\Gamma \Proves !A$ و$\lnot !A \in \Gamma$، فإن $\Gamma$
  غير متسقة.
\end{prop}

\begin{proof}
  افترض أن $\Gamma \Proves !A$ و$\lnot !A \in \Gamma$. إذن توجد
  $!B_1$, \dots, $!B_n \in \Gamma$ بحيث
  \[
  \{\sFmla{\False}{!A}, \sFmla{\True}{!B_1}, \dots, \sFmla{\True}{!B_n}\}
  \]
  لها جدول دلالي مغلق. استبدل بالافتراض \sFmla{\False}{!A} الافتراض
  \sFmla{\True}{\lnot !A}، وأدرج بعد الافتراضات نتيجة تطبيق
  \TRule{\True}{\lnot} على \sFmla{\True}{\lnot !A}. وكل !!{sentence}
  في ال!!{tableau} كان تبريرها يحيل إلى السطر~$1$ في ال!!{tableau}
  القديم يصبح تبريرها الآن
  محيلًا إلى السطر~$n+2$. ولذلك إذا كان ال!!{tableau} القديم مغلقًا، كان
  الجديد مغلقًا أيضًا. وهو يثبت أن $\Gamma$ غير متسقة، لأن جميع الافتراضات
  تنتمي إلى~$\Gamma$.
\end{proof}

\begin{prop}\ollabel{prop:provability-exhaustive}
  إذا كانت كل من $\Gamma \cup \{!A\}$ و$\Gamma \cup \{\lnot !A\}$
  غير متسقة، فإن $\Gamma$ غير متسقة.
\end{prop}

\begin{proof}
  إذا وُجدت $!B_1$, \dots, $!B_n \in \Gamma$ و$!C_1$, \dots,
  $!C_m \in \Gamma$ بحيث
  \begin{align*}
    \{\sFmla{\True}{!A}, &
    \sFmla{\True}{!B_1}, \dots, \sFmla{\True}{!B_n}\} \text{ و}\\
    \{\sFmla{\True}{\lnot !A}, &
    \sFmla{\True}{!C_1}, \dots, \sFmla{\True}{!C_m}\}
  \end{align*}
  كان لكل منهما !!{tableau}s مغلقة، فيمكننا إنشاء !!{tableau} واحد مركب
  يثبت أن $\Gamma$ غير متسقة، وذلك باستخدام $\sFmla{\True}{!B_1}$،
  \dots، $\sFmla{\True}{!B_n}$ مع $\sFmla{\True}{!C_1}$، \dots،
  $\sFmla{\True}{!C_m}$ بوصفها افتراضات، ثم تطبيق قاعدة \Cut{}. وينتج
  فرعان، يبدأ أحدهما بــ$\sFmla{\True}{!A}$ والآخر بــ
  $\sFmla{\False}{!A}$.  
  
  في الجانب الأيسر، أضف الجزء الواقع تحت افتراضات ال!!{tableau} الأول.
  وهنا يظل كل تطبيق لقاعدة صحيحًا، لأن كل افتراض من افتراضات ال!!{tableau}
  الأول، بما فيها $\sFmla{\True}{!A}$، متاح. وهكذا يُغلق كل فرع تحت
  $\sFmla{\True}{!A}$. 
  
  وفي الجانب الأيمن، أضف الجزء الواقع تحت افتراضات ال!!{tableau} الثاني،
  بعد حذف نتائج أي تطبيقات لــ~$\TRule{\True}{\lnot}$ على
  $\sFmla{\True}{\lnot !A}$. فنتيجة $\TRule{\True}{\lnot}$ عند تطبيقها
  على $\sFmla{\True}{\lnot !A}$ هي $\sFmla{\False}{!A}$، وهي مع ذلك
  متاحة لأنها نتيجة قاعدة \Cut{} في الجانب الأيمن من ال!!{tableau} المركب.

  إذا كان فرع في الجدول الدلالي الثاني مغلقًا لاحتوائه على الافتراض
  $\sFmla{\True}{\lnot !A}$ (الذي لم يعد يظهر افتراضًا في ال!!{tableau}
  المركب) وعلى $\sFmla{\False}{\lnot !A}$، فيمكننا تطبيق
  $\TRule{\False}{\lnot}$ على $\sFmla{\False}{\lnot !A}$ لنحصل على
  $\sFmla{\True}{!A}$. وعندئذ يُغلق أيضًا الفرع الموافق في ال!!{tableau}
  المركب، لأنه يحتوي على نتيجة قاعدة \Cut{} اليمنى،
  $\sFmla{\False}{!A}$. وإذا أُغلق فرع في ال!!{tableau} الثاني لأي سبب
  آخر، فإن الفرع الموافق له في ال!!{tableau} المركب يُغلق أيضًا، لأن كل
  !!{signed formula}s غير $\sFmla{\True}{\lnot !A}$ تظهر في الفرع من
  ال!!{tableau} الثاني القديم تظهر كذلك في الفرع الموافق من ال!!{tableau}
  المركب.
  \end{proof}

\end{document}
